%% FullCourse_10Lecs -- Lecture 10, Topic 10.1: Framing
%% Source: ./archive_pre_split/Lec10_Case_Studies_v2.tex (section 1)

%% Shared preamble for the FullCourse_10Lecs topic decks.
%%
%% Each topic .tex \input's this file, then defines its own \title, \date
%% override (if any), \graphicspath, and \begin{document}.
%%
%% Reconstructed verbatim from the inlined copy preserved in
%% Lec10_Case_Studies/Lec10_T8_Case6_DeepRamsey.tex (lines 23-190), which is
%% explicitly annotated there as "from ../shared_preamble.tex, copied verbatim".
%% Base preamble matches MiniCourse_8hr/Slides/shared_preamble.tex; this version
%% additionally provides \sectiondivider, the conceptbox/cautionbox/examplebox/
%% takeawaybox tcolorboxes, and \compacttables used across the split decks.

\documentclass[10pt,english,aspectratio=169]{beamer}

\usepackage{ae,aecompl}
\usepackage[T1]{fontenc}
\usepackage{textcomp}
\usepackage[utf8]{inputenc}
\usepackage{babel}
\usepackage{datetime2}

\setcounter{secnumdepth}{3}
\setcounter{tocdepth}{3}

\usepackage{amsmath, amssymb, amsthm, graphicx}
\usepackage{booktabs, multirow}
\usepackage{tikz}
\usepackage{pgfplots}
\pgfplotsset{compat=1.16}
\usepackage[most]{tcolorbox}
\tcbuselibrary{skins, breakable, theorems}
\usepackage{listings}
\usepackage{xcolor}

\lstset{
  basicstyle=\ttfamily\small,
  keywordstyle=\color{blue!80!black}\bfseries,
  commentstyle=\color{green!50!black}\itshape,
  stringstyle=\color{orange!80!black},
  backgroundcolor=\color{gray!8},
  frame=single,
  rulecolor=\color{gray!40},
  breaklines=true,
  showstringspaces=false,
  numbers=none,
}

\lstdefinestyle{pythonstyle}{
    language=Python,
    basicstyle=\ttfamily\footnotesize,
    keywordstyle=\color{blue!80!black}\bfseries,
    commentstyle=\color{green!50!black}\itshape,
    stringstyle=\color{orange!80!black},
    frame=single,
    breaklines=true,
    showstringspaces=false,
    numbers=left,
    numbersep=5pt,
    numberstyle=\tiny\color{gray},
    backgroundcolor=\color{gray!8},
    rulecolor=\color{gray!40}
}

\ifx\hypersetup\undefined
  \AtBeginDocument{%
    \hypersetup{bookmarks=false,breaklinks=true,pdfborder={0 0 0},colorlinks=true,
      linkcolor=DarkText,urlcolor=RedTitle}%
  }
\else
  \hypersetup{bookmarks=false,breaklinks=true,pdfborder={0 0 0},colorlinks=true,
    linkcolor=DarkText,urlcolor=RedTitle}%
\fi

\makeatletter
\newcommand\makebeamertitle{%
  \begin{frame}[plain]
    \vspace*{1.0cm}
    \centering
    {\usebeamerfont{title}\usebeamercolor[fg]{title}\inserttitle\par}
    \vspace{1.0cm}
    {\usebeamerfont{author}\usebeamercolor[fg]{author}\insertauthor\par}
    \vspace{0.2cm}
    {\usebeamerfont{date}\usebeamercolor[fg]{date}\insertdate\par}
  \end{frame}}%
\AtBeginDocument{%
  \let\origtableofcontents=\tableofcontents
  \def\tableofcontents{\@ifnextchar[{\origtableofcontents}{\gobbletableofcontents}}
  \def\gobbletableofcontents#1{\origtableofcontents}%
}

\usetheme{metropolis}
\usefonttheme{professionalfonts}

\definecolor{LightGray}{RGB}{230,230,230}
\definecolor{RedTitle}{RGB}{0,0,200}
\definecolor{VeryLightGray}{RGB}{245,245,245}
\definecolor{DarkText}{RGB}{0,0,0}

\setbeamercolor{background canvas}{bg=white}
\setbeamercolor{title}{fg=RedTitle}
\setbeamercolor{author}{fg=DarkText}
\setbeamercolor{date}{fg=DarkText}
\setbeamercolor{frametitle}{bg=VeryLightGray,fg=DarkText}
\setbeamercolor{normal text}{fg=DarkText}
\setbeamerfont{title}{size=\large}

\setbeamersize{text margin left=5mm, text margin right=5mm}
\addtobeamertemplate{frametitle}{}{\vspace{-0.5em}}
\newcommand{\reducevspace}{\vspace{-0.5cm}}

% Shared section divider and box styles for split topic decks.
\newcommand{\sectiondivider}[2]{%
\begin{frame}[plain,noframenumbering]
\begin{center}
\vfill
{\Large\textbf{#1}\par}
\vspace{0.35cm}
{\normalsize #2\par}
\vfill
\end{center}
\end{frame}
}

\newtcolorbox{conceptbox}[1][]{
  colback=blue!3,
  colframe=RedTitle!55,
  boxrule=0.35mm,
  arc=1mm,
  left=2mm,
  right=2mm,
  top=1mm,
  bottom=1mm,
  #1
}
\newtcolorbox{cautionbox}[1][]{
  colback=red!3,
  colframe=red!50!black,
  boxrule=0.35mm,
  arc=1mm,
  left=2mm,
  right=2mm,
  top=1mm,
  bottom=1mm,
  #1
}
\newtcolorbox{examplebox}[1][]{
  colback=green!3,
  colframe=green!45!black,
  boxrule=0.35mm,
  arc=1mm,
  left=2mm,
  right=2mm,
  top=1mm,
  bottom=1mm,
  #1
}
\newtcolorbox{takeawaybox}[1][]{
  colback=VeryLightGray,
  colframe=RedTitle!55,
  boxrule=0.35mm,
  arc=1mm,
  left=2mm,
  right=2mm,
  top=1mm,
  bottom=1mm,
  #1
}
\newcommand{\compacttables}{\renewcommand{\arraystretch}{1.15}}

\setbeamertemplate{navigation symbols}{}
\setbeamertemplate{footline}{%
  \hfill%
  \usebeamercolor[fg]{page number in head/foot}%
  \usebeamerfont{page number in head/foot}%
  \insertframenumber\,/\,\inserttotalframenumber\kern1em\vskip2pt%
}

\makeatother

\author{Zhigang Feng}
% "date with time": compile timestamp so each build is identifiable.
\date{\today\ \DTMcurrenttime}


% Suppress redundant auto-generated metropolis section page; \sectiondivider frames are the dividers we keep.
\metroset{sectionpage=none}

\graphicspath{{./pic/}{./figures/}{./Exercise10_ES_Fellows/plots/}{../../MiniCourse_8hr/Slides/Module4_Agentic_AI/pic/}{../}}

\usetikzlibrary{positioning,arrows.meta,shapes,calc,fit,backgrounds}

\lstdefinestyle{markdownstyle}{
    basicstyle=\ttfamily\footnotesize,
    frame=single,
    rulecolor=\color{gray!40},
    backgroundcolor=\color{gray!8},
    breaklines=true,
    showstringspaces=false,
    numbers=none,
}

\title[AI for Econ Research]{AI for Economic Research: Dynamic Models, Language, and Agents\\
Lecture 10.1: Case Studies --- Framing}

\begin{document}

\makebeamertitle

%=============================================================================
\begin{frame}{Topic 10.1 Agenda}
\begin{enumerate}
  \item \textbf{Framing} --- why case studies, what they show, road map for the lecture
\end{enumerate}
\end{frame}

%=============================================================================
\section{Framing}
%===========================================================

\sectiondivider{Framing}{From Lec.~9 concepts to complete research workflows}

\begin{frame}{Where We Are in the Course}
\begin{center}
\begin{tikzpicture}[
    node distance=0.25cm,
    lecture/.style={rectangle, rounded corners, draw=gray!60, fill=gray!10,
        minimum width=3.5cm, minimum height=0.45cm, font=\footnotesize, text=DarkText},
    current/.style={rectangle, rounded corners, draw=RedTitle, fill=blue!8,
        minimum width=3.5cm, minimum height=0.45cm, font=\footnotesize\bfseries, text=RedTitle},
]
\node[lecture] (l1) {1--6: models + AI/ML tools};
\node[lecture, right=0.3cm of l1] (l2) {7--8: LLMs + RAG};
\node[lecture, right=0.3cm of l2] (l3) {9: agentic AI workflow};
\node[current, below=0.25cm of l2] (l4) {10: case studies};
\end{tikzpicture}
\end{center}

\vspace{0.3cm}
\begin{itemize}
    \item Lec.~9 gave the workflow: specification, structure, validation, and human oversight
    \item Lec.~10 shows how that workflow changes shape across different research jobs
    \item The core question is no longer ``what is a skill?'' It is ``what workflow should I build for this task?''
\end{itemize}
\end{frame}

\begin{frame}{Lab Readiness Check}

Before today's exercises, verify all five:

\vspace{0.2cm}

\begin{enumerate}
    \item \textbf{Terminal agent ready} --- one of Claude Code, Codex, Gemini CLI, Cursor, or a course-approved equivalent can run

    \vspace{0.1cm}

    \item \textbf{Access confirmed} --- login, quota, and model access checked before class

    \vspace{0.1cm}

    \item \textbf{Git initialized in your project} --- \texttt{git status} runs without error

    \vspace{0.1cm}

    \item \textbf{Working Python or R environment} --- \texttt{which python} or \texttt{which R} returns a path

    \vspace{0.1cm}

    \item \textbf{Project memory created} --- \texttt{AGENTS.md}, \texttt{CLAUDE.md}, \texttt{CODEX.md}, or equivalent rules file with project name, language, and goal
\end{enumerate}

\vspace{0.3cm}

\textbf{If any box is unchecked,} open Appendix~A now. The exercises in this lecture require all five. If setup takes more than 10 minutes, pair with a neighbor and share one working project.
\end{frame}

\begin{frame}{Why Case Studies After Lec.~9}
\begin{itemize}
    \item Knowing the vocabulary of agentic AI does not tell you how to design a research workflow
    \item The difficult step is decomposition: where to split the task, where to validate, and what stays human
    \item Different tasks need different harnesses:
    paper review is not the same as MEPS harmonization, AI-exposure measurement, OLG homotopy, or public-web metadata collection
    \item Today is about design judgment, not product loyalty
\end{itemize}

\vspace{0.35cm}
\textbf{Question for every case:} what is the artifact, what is the failure mode, and what must the economist still verify?
\end{frame}

\begin{frame}{Three Research Patterns}
\small
\begin{tabular}{p{3.0cm}p{4.1cm}p{5.2cm}}
\toprule
\textbf{Pattern} & \textbf{Representative task} & \textbf{Main risk if you use the wrong workflow} \\
\midrule
Reusable skill & Referee-style paper review & Generic comments, hallucinated quotes, missing quality gate \\
Empirical pipeline & MEPS harmonization and figures & Wrong crosswalk, dropped survey weights, silent validation failure \\
Public-web dataset & Econometric Society fellows metadata & Broken source discipline, missing provenance, hidden missingness \\
\bottomrule
\end{tabular}

\vspace{0.3cm}
\textbf{All three use agents, but the harness is different.}
\end{frame}

\begin{frame}{What Remains Human Across All Three}
\begin{columns}[T]
\begin{column}{0.48\textwidth}
\textbf{You still own}
\begin{itemize}
    \item the research question
    \item the coding rules
    \item the validation target
    \item the interpretation
    \item the decision to trust or reject the output
\end{itemize}
\end{column}
\begin{column}{0.48\textwidth}
\textbf{The agent accelerates}
\begin{itemize}
    \item repetitive execution
    \item file handling
    \item promptable decomposition
    \item re-runs after fixes
    \item draft diagnostics and plots
\end{itemize}
\end{column}
\end{columns}

\vspace{0.35cm}
\textbf{Lec.~9 still governs this lecture:} specification and validation are not optional just because the execution is faster.
\end{frame}

\end{document}
